% !TeX root = ../main.tex
\section{Discussion}
\label{sec:discussion}

% --------------------------------------------------------
\subsection{Adaptive detection of recurrent plantation harvests}
\label{sec:discussion_local_adaptation}

Recurrent harvest detection depends on changes relative to evolving canopy conditions. Although pooled pre- and post-harvest backscatter distributions overlapped, within-event trajectories showed a consistent decline followed by recovery. Temporal trajectories therefore provide information that static backscatter thresholds cannot capture. Similar post-harvest responses have been reported in other short-rotation plantations~\citep{Spracklen2021Acacia}. VV also supported better event detection than VH despite comparable median declines, indicating that harvest contrast alone does not determine detection performance. Background variability during regrowth may also affect the separation of harvest signals from other fluctuations~\citep{Sabir2024AdaptiveCuSUM}.

Over successive rotations, a full-record baseline combines different stages of stand development~\citep{Kennedy2010LandTrendr}. Windowed CuSum renews this baseline locally, while central decision intervals and permutation testing constrain candidate acceptance. Unlike recursive ReCuSum~\citep{Ygorra2023ReCuSum}, later searches do not depend on earlier segmentation decisions. This distinction matters when an early harvest is missed or misdated. The retrospective design also uses temporal context beyond the candidate event, unlike methods intended for immediate disturbance alerts~\citep{Ygorra2024CuSumNRT}.

The comparison with first-post ReCuSum shows that the main gain was more complete harvest histories, with similar dating accuracy on shared detections. Recovering successive events on the same land is essential for estimating harvest frequency and recurrence intervals; accurately dating isolated events alone cannot provide this information.

% --------------------------------------------------------
\subsection{Interpreting balsa harvest dynamics, hotspots, and recurrence intervals}
\label{sec:discussion_harvest_histories}

Harvest operations can span several months and spatial sub-blocks, creating a mismatch between pixelwise changes and management events. TDA addresses this mismatch by consolidating detections while distinguishing observed inactivity from insufficient observation support. STP adds spatial detail within those events. Together, they support interpretation at both the stand and sub-stand scales without treating each harvesting stage as another rotation.

The national reconstruction adds a temporal dimension to the plantation inventory, identifying where harvesting recurs and when activity concentrates. The western concentration and the 2022 peak followed by a sharp decline reveal substantial spatial and temporal variation in canopy turnover. These patterns can guide comparisons with timber production and planting records; the satellite histories alone do not establish market demand as their cause.

The national median interval of 34 months is consistent with balsa reaching commercial dimensions within approximately three years on productive sites~\citep{GonzalezOsorio2010Balsa,CanadasLopez2019BalsaGrowth}. However, the dominance of Am plantations means that this national estimate primarily describes that climate zone. The longer modal interval in Aw plantations may reflect slower growth under seasonal water limitation~\citep{Yuan2026BalsaWaterUse}, although management differences could also contribute. The 3--5-month mode in the smaller Af sample cannot represent a complete biological rotation; possible explanations include staged clearing and non-harvest backscatter changes. These intervals describe detected stand-level canopy turnover and need not equal planting-to-harvest rotation lengths.

% --------------------------------------------------------
\subsection{Limitations, transferability, and future perspectives}
\label{sec:discussion_limitations}

C-band backscatter changes do not uniquely identify clear-felling. Rainfall, flooding, windthrow, and thinning can produce declines~\citep{Doblas2020Rainfall}, while residual trees or slash can weaken harvest signals~\citep{Balling2023Texture}. Optical reference records may also omit cloud-obscured events, so unmatched detections can include undocumented harvests as well as false detections. Field or production records would help distinguish these outcomes.

Spatial support and temporal assumptions further limit interpretation. The 30~m grid and stand-support threshold can omit small or narrow cuts. Permutation testing assumes exchangeable observations; the sharp loss of recall under calendar-block permutations shows that temporal dependence materially affects detection. This sensitivity requires further evaluation rather than treating the nominal threshold as a guaranteed false-alarm rate. Dependence on post-event observations also makes the framework most suitable for retrospective reconstruction, with reduced temporal support near the archive endpoint.

Transfer to other short-rotation plantations, including \textit{Eucalyptus} and \textit{Acacia}~\citep{Spracklen2021Acacia,Huang2026EucalyptusAge}, would require evaluating harvest responses and adapting temporal windows to local growth and management. Future work could combine complementary SAR frequencies with independent harvest records to distinguish canopy removal from moisture-driven changes and assess recovery across successive rotations.